Advancements in semiconductor technology have historically driven computational power growth. For decades, the number of transistors on a chip doubled roughly every eighteen months, following the observation formalized as \textit{Moore's Law} \cite{Pacheco2011}; rising clock speeds translated these denser circuits into proportional single-core performance gains. However, in the mid-2000s, the hardware industry encountered significant physical barriers, such as excessive power consumption and challenges in heat dissipation, which limited the increase in processor clock frequency.

The response to these technological limitations was a change in processor design: instead of focusing on making a single core faster, the emphasis shifted to integrating multiple processing cores onto a single chip. Thus, multi-core processors emerged as the dominant architecture in modern computers. This transition opened up new possibilities for performance enhancement; however, to take advantage of the performance gains offered by multi-core processors, programs must be intentionally developed or adapted to execute tasks in parallel \cite{Pacheco2011}.

This architectural shift from sequential to parallel execution profoundly impacts high-performance applications. In many scenarios, such as bioinformatics pipelines and network intrusion detection, processing large data volumes efficiently is essential. At the heart of these applications lies the need for highly efficient multi-pattern-matching algorithms, among which the Aho-Corasick algorithm stands out for its efficiency. Widely adopted in security tools --- including YARA for malware classification and Snort and Suricata for signature-based network intrusion detection \cite{LeeYang2017FHBM} --- the Aho-Corasick algorithm serves as a backbone for multi-pattern search engines \cite{YARADocs}.

Despite its efficiency, the Aho-Corasick algorithm is intrinsically sequential: its automaton processes input strictly left-to-right, with each state transition depending on the preceding one. As a result, when scanning a single large input on a modern multi-core CPU, the entire computational load is concentrated on one core. This serial dependency becomes a significant performance bottleneck as the size of the target data increases, such as with full disk images, packet captures, or genomic sequences. Consequently, the disparity between the algorithm's sequential nature and modern hardware's parallel processing capabilities widens as data volumes grow \cite{AhoCorasick1975}.

\newpage\subsecao{Motivation}

The intrinsically sequential nature of Aho-Corasick, combined with the prevalence of multi-core architectures in general-purpose computing, presents a clear opportunity for acceleration through parallel computing. Prior work reports speedups on specialized platforms and in specific application domains \cite{Vajira2018, LeeYang2017FHBM}; however, the behavior of a shared-memory, multi-threaded implementation on commodity CPUs in the \textit{memory-bound regime} --- where the automaton size far exceeds the last-level cache --- has received comparatively little attention, particularly on modern hybrid processors that combine high-performance and high-efficiency cores.

As pattern sets for tools such as Snort and Suricata grow in size, the corresponding Aho-Corasick automata can occupy hundreds of megabytes, pushing the hot data structure well beyond the cache hierarchy \cite{Aldwairi2018Characterizing}. In this regime, the dominant performance limit shifts from computation to memory bandwidth, and the interaction of multiple threads sharing the same last-level cache changes the scaling behavior in ways that compute-bound models do not capture. Understanding these limits on realistic workloads is therefore a prerequisite for deploying parallelized pattern-matching engines in production NIDS pipelines.

This thesis addresses that gap through the following research questions:
\begin{itemize}
    \item \textbf{RQ1:} What throughput and speedup gains does text-level data parallelism, implemented with POSIX Threads over a shared read-only automaton, yield for Aho-Corasick on a shared-memory multi-core CPU across different workload and pattern-set configurations?
    \item \textbf{RQ2:} How does the automaton-to-last-level-cache size ratio constrain parallel scalability, and what is the dominant performance ceiling in the memory-bound regime on a hybrid P/E-core processor?
\end{itemize}

\subsecao{Objectives}

The general objective of this thesis is to design, implement, and evaluate a parallel version of the Aho-Corasick algorithm targeting shared-memory multi-core CPU architectures, using the POSIX Threads programming model, in order to accelerate multi-pattern matching on large inputs.

To achieve this general objective, the following specific objectives are defined:
\begin{enumerate}
    \item To design and implement a parallel version of the Aho-Corasick algorithm in C, using the pthreads API, exploiting intra-input parallelism on shared-memory multi-core CPUs;
    \item To conduct a rigorous experimental evaluation that quantifies the performance gains of the parallel implementation in terms of speedup, throughput, and parallel efficiency, against a sequential baseline;
    \item To analyze the scalability of the proposed solution under different workload, pattern-set, and thread-count configurations.
\end{enumerate}

\subsecao{Structure of the Thesis}

This document is organized as follows. Chapter 2 presents the Theoretical Framework, introducing the fundamentals of parallel computing, including architectures, performance metrics, and thread-based programming models, with emphasis on the POSIX Threads API. It also discusses the core data structures and algorithms relevant to this work, including Tries, the general pattern-matching problem, and a detailed analysis of the Aho-Corasick algorithm. Chapter 3 describes the Systematic Literature Review, presenting the research questions, search strategy, and selection criteria adopted to identify the state of the art and define the research gap. Chapter 4 presents the Proposed Solution, detailing the parallel-scanning strategy, the boundary-overlap and load-balancing mechanisms, and the flattened output layout that constitute the design. Chapter 5 describes the Methodology, specifying the experimental environment, datasets, metrics, measurement protocol, correctness validation, and the hypotheses that the experiments are designed to test. Chapter 6 reports and discusses the Results, isolating the contribution of each optimization and positioning the findings against the related literature. Finally, Chapter 7 presents the Conclusion, summarizing the contributions, the threats to validity, and directions for future work.
% Fim do conteúdo da seção de Introdução.

